% Options for packages loaded elsewhere
\PassOptionsToPackage{unicode}{hyperref}
\PassOptionsToPackage{hyphens}{url}
\documentclass[
  nobib]{tufte-handout}
\usepackage{xcolor}
\usepackage{amsmath,amssymb}
\setcounter{secnumdepth}{-\maxdimen} % remove section numbering
\usepackage{iftex}
\ifPDFTeX
  \usepackage[T1]{fontenc}
  \usepackage[utf8]{inputenc}
  \usepackage{textcomp} % provide euro and other symbols
\else % if luatex or xetex
  \usepackage{unicode-math} % this also loads fontspec
  \defaultfontfeatures{Scale=MatchLowercase}
  \defaultfontfeatures[\rmfamily]{Ligatures=TeX,Scale=1}
\fi
\usepackage{lmodern}
\ifPDFTeX\else
  % xetex/luatex font selection
\fi
% Use upquote if available, for straight quotes in verbatim environments
\IfFileExists{upquote.sty}{\usepackage{upquote}}{}
\IfFileExists{microtype.sty}{% use microtype if available
  \usepackage[]{microtype}
  \UseMicrotypeSet[protrusion]{basicmath} % disable protrusion for tt fonts
}{}
\makeatletter
\@ifundefined{KOMAClassName}{% if non-KOMA class
  \IfFileExists{parskip.sty}{%
    \usepackage{parskip}
  }{% else
    \setlength{\parindent}{0pt}
    \setlength{\parskip}{6pt plus 2pt minus 1pt}}
}{% if KOMA class
  \KOMAoptions{parskip=half}}
\makeatother
\usepackage{color}
\usepackage{fancyvrb}
\newcommand{\VerbBar}{|}
\newcommand{\VERB}{\Verb[commandchars=\\\{\}]}
\DefineVerbatimEnvironment{Highlighting}{Verbatim}{commandchars=\\\{\}}
% Add ',fontsize=\small' for more characters per line
\newenvironment{Shaded}{}{}
\newcommand{\AlertTok}[1]{\textcolor[rgb]{1.00,0.00,0.00}{\textbf{#1}}}
\newcommand{\AnnotationTok}[1]{\textcolor[rgb]{0.38,0.63,0.69}{\textbf{\textit{#1}}}}
\newcommand{\AttributeTok}[1]{\textcolor[rgb]{0.49,0.56,0.16}{#1}}
\newcommand{\BaseNTok}[1]{\textcolor[rgb]{0.25,0.63,0.44}{#1}}
\newcommand{\BuiltInTok}[1]{\textcolor[rgb]{0.00,0.50,0.00}{#1}}
\newcommand{\CharTok}[1]{\textcolor[rgb]{0.25,0.44,0.63}{#1}}
\newcommand{\CommentTok}[1]{\textcolor[rgb]{0.38,0.63,0.69}{\textit{#1}}}
\newcommand{\CommentVarTok}[1]{\textcolor[rgb]{0.38,0.63,0.69}{\textbf{\textit{#1}}}}
\newcommand{\ConstantTok}[1]{\textcolor[rgb]{0.53,0.00,0.00}{#1}}
\newcommand{\ControlFlowTok}[1]{\textcolor[rgb]{0.00,0.44,0.13}{\textbf{#1}}}
\newcommand{\DataTypeTok}[1]{\textcolor[rgb]{0.56,0.13,0.00}{#1}}
\newcommand{\DecValTok}[1]{\textcolor[rgb]{0.25,0.63,0.44}{#1}}
\newcommand{\DocumentationTok}[1]{\textcolor[rgb]{0.73,0.13,0.13}{\textit{#1}}}
\newcommand{\ErrorTok}[1]{\textcolor[rgb]{1.00,0.00,0.00}{\textbf{#1}}}
\newcommand{\ExtensionTok}[1]{#1}
\newcommand{\FloatTok}[1]{\textcolor[rgb]{0.25,0.63,0.44}{#1}}
\newcommand{\FunctionTok}[1]{\textcolor[rgb]{0.02,0.16,0.49}{#1}}
\newcommand{\ImportTok}[1]{\textcolor[rgb]{0.00,0.50,0.00}{\textbf{#1}}}
\newcommand{\InformationTok}[1]{\textcolor[rgb]{0.38,0.63,0.69}{\textbf{\textit{#1}}}}
\newcommand{\KeywordTok}[1]{\textcolor[rgb]{0.00,0.44,0.13}{\textbf{#1}}}
\newcommand{\NormalTok}[1]{#1}
\newcommand{\OperatorTok}[1]{\textcolor[rgb]{0.40,0.40,0.40}{#1}}
\newcommand{\OtherTok}[1]{\textcolor[rgb]{0.00,0.44,0.13}{#1}}
\newcommand{\PreprocessorTok}[1]{\textcolor[rgb]{0.74,0.48,0.00}{#1}}
\newcommand{\RegionMarkerTok}[1]{#1}
\newcommand{\SpecialCharTok}[1]{\textcolor[rgb]{0.25,0.44,0.63}{#1}}
\newcommand{\SpecialStringTok}[1]{\textcolor[rgb]{0.73,0.40,0.53}{#1}}
\newcommand{\StringTok}[1]{\textcolor[rgb]{0.25,0.44,0.63}{#1}}
\newcommand{\VariableTok}[1]{\textcolor[rgb]{0.10,0.09,0.49}{#1}}
\newcommand{\VerbatimStringTok}[1]{\textcolor[rgb]{0.25,0.44,0.63}{#1}}
\newcommand{\WarningTok}[1]{\textcolor[rgb]{0.38,0.63,0.69}{\textbf{\textit{#1}}}}
\usepackage{longtable,booktabs,array}
\newcounter{none} % for unnumbered tables
\usepackage{calc} % for calculating minipage widths
% Correct order of tables after \paragraph or \subparagraph
\usepackage{etoolbox}
\makeatletter
\patchcmd\longtable{\par}{\if@noskipsec\mbox{}\fi\par}{}{}
\makeatother
% Allow footnotes in longtable head/foot
\IfFileExists{footnotehyper.sty}{\usepackage{footnotehyper}}{\usepackage{footnote}}
\makesavenoteenv{longtable}
\setlength{\emergencystretch}{3em} % prevent overfull lines
\providecommand{\tightlist}{%
  \setlength{\itemsep}{0pt}\setlength{\parskip}{0pt}}
% =====================================================================
% axona-doc-preamble.tex — shared LaTeX preamble for the programmer-guide
% quartet (Quick-Start, Programmer Guide, API Reference, Services Guide).
%
% These docs are markdown-native; the PDF is built by converting the .md to
% a tufte-handout .tex with pandoc (-H this file) and compiling with tectonic.
% See render.sh. Modeled on explainer/Axona-Explainer.tex (the flat-essay
% original); the tweaks below make tufte-handout work for code/table-heavy
% reference material.
% =====================================================================

\definecolor{rust}{RGB}{185,78,54}
\definecolor{rustsoft}{RGB}{212,168,150}
\definecolor{inkmuted}{RGB}{102,102,102}
\usepackage{microtype}

% Explainer-matched page furniture: tufte-handout's letterspaced running
% heads fail on certain page breaks (the SOUL \MakeTextLowercase issue the
% explainer works around), so use the same fancyhdr style — no running
% head, a centered page number in the footer.
\usepackage{fancyhdr}
\fancypagestyle{ndhtplain}{%
  \fancyhf{}%
  \renewcommand{\headrulewidth}{0pt}%
  \fancyfoot[C]{\thepage}%
}

% Explainer-matched heading-orphan protection: reserve vertical room before
% each (sub)section heading so a heading never sits alone at a page bottom.
\usepackage{etoolbox}
\usepackage{needspace}
\pretocmd{\section}{\needspace{6\baselineskip}}{}{}
\pretocmd{\subsection}{\needspace{4\baselineskip}}{}{}

% XeTeX (tectonic) unicode monospace: renders … → and the box-drawing ASCII
% diagrams that the default CM/LM mono cannot. Menlo ships with macOS.
\setmonofont{Menlo}[Scale=0.78]

% tufte-handout ships \subsubsection as an error stub (it wants flat essays).
% A run-in fallback lets nested reference headings (§4.2, §12.3, …) compile.
\renewcommand{\subsubsection}[1]{\medskip\par\noindent{\normalsize\bfseries #1}\par\nobreak\smallskip}

% These docs use no margin notes, so reclaim tufte's wide right margin: widen
% the text block so code blocks and API tables don't overflow it.
\usepackage{geometry}
\geometry{left=1in, textwidth=6.1in, marginparsep=0.2in, marginparwidth=0.9in,
          top=1in, headsep=0.4in}

% Wrap long code lines (cheat-sheet comments run to ~130 chars) with a ↪ marker
% instead of overflowing. Applies to pandoc's fancyvrb Highlighting/Verbatim.
\usepackage{fvextra}
\fvset{breaklines=true, breakanywhere=true}

% A few technical symbols appear in prose / inline code that Latin Modern and
% Menlo don't cover; map them to math equivalents so they render (rather than
% dropping to tofu). Only maps chars actually used — leaves → alone (it renders).
\usepackage{newunicodechar}
\newunicodechar{≈}{\ensuremath{\approx}}
\newunicodechar{≠}{\ensuremath{\neq}}
\newunicodechar{≥}{\ensuremath{\geq}}
\newunicodechar{≤}{\ensuremath{\leq}}
\newunicodechar{×}{\ensuremath{\times}}
\newunicodechar{·}{\ensuremath{\cdot}}
\newunicodechar{‖}{\ensuremath{\|}}
\usepackage{bookmark}
\IfFileExists{xurl.sty}{\usepackage{xurl}}{} % add URL line breaks if available
\urlstyle{same}
\hypersetup{
  hidelinks,
  pdfcreator={LaTeX via pandoc}}

\author{}
\date{}

\begin{document}

% Generated by render.sh — the explainer-style title page.
\thispagestyle{empty}
\null\vfill
\begin{center}
{\fontsize{48}{52}\selectfont \textsc{Axona}\par}
\vspace{0.6em}
{\Large\itshape A Working Pub/Sub Roundtrip in Five Minutes\par}
\vspace{4em}
{\large David A.~Smith\par}
{\itshape\small Axona.net\par}
\vspace{2em}
{\small The Axona Quick Start \quad v4.40.0 \quad 2026-07-24\par}
\end{center}
\vfill
\null
\newpage
\pagestyle{ndhtplain}
\thispagestyle{ndhtplain}
\setcounter{page}{1}

\section{Axona Quick Start}\label{axona-quick-start}

Get a working pub/sub roundtrip in \textbf{under five minutes} on the
current \texttt{@axona/protocol} \textbf{v4.40.0} API (kernel 4.40.0).
One Node process connects to the live \textbf{testnet} bridge,
subscribes to an open topic, publishes a signed message, and logs what
comes back.

\begin{quote}
\textbf{Two live networks, one line.} Both run the 4.x kernel:
\textbf{testnet} (\texttt{wss://testnet.axona.net}) tracks the newest
release --- 4.40.0, which this Quick Start pins --- and
\textbf{production} (\texttt{wss://bridge.axona.net}) runs the most
recently promoted release. This Quick Start uses testnet; swap the
bridge URL to target production.
\end{quote}

Companion documents:

\begin{itemize}
\tightlist
\item
  \href{Axona-Programmer-Guide-v4.38.0.md}{Programmer Guide} --- the
  five ideas + recipes for real apps.
\item
  \href{Axona-API-Reference-v4.40.0.md}{API Reference} --- every
  exported symbol.
\item
  \href{Axona-AI-Grounding-v4.40.0.md}{AI Grounding} --- building with
  an AI assistant? Hand it this file.
\end{itemize}

\subsection{Prerequisites}\label{prerequisites}

\begin{itemize}
\tightlist
\item
  \textbf{Node.js 20+} (for built-in Web Crypto Ed25519)
\item
  A terminal with network access (we connect to
  \texttt{wss://testnet.axona.net})
\end{itemize}

No build step, no DB, no local bridge --- the testnet bridge is the
entry point.

\subsection{1. Install (30 seconds)}\label{install-30-seconds}

\begin{Shaded}
\begin{Highlighting}[]
\FunctionTok{mkdir}\NormalTok{ my{-}axona{-}demo }\KeywordTok{\&\&} \BuiltInTok{cd}\NormalTok{ my{-}axona{-}demo}
\ExtensionTok{npm}\NormalTok{ init }\AttributeTok{{-}y}
\ExtensionTok{npm}\NormalTok{ pkg set type=module}
\ExtensionTok{npm}\NormalTok{ install github:axona{-}net/axona{-}protocol\#v4.40.0}
\end{Highlighting}
\end{Shaded}

\subsection{2. Two identities, one rule}\label{two-identities-one-rule}

Axona separates \textbf{who you connect as} from \textbf{who you publish
as}:

\begin{itemize}
\tightlist
\item
  \texttt{createNodeIdentity(\{\ lat,\ lng\ \})} -\textgreater{} your
  \textbf{connection} identity. Its public key forms the 66-hex node ID
  (\texttt{.id}); the leading byte encodes your region. This
  authenticates the transport. It is fine to mint a fresh one each run.
\item
  \texttt{createAuthorIdentity(\{\ persistAs?\ \})} -\textgreater{} your
  \textbf{author} identity. Its public key is the 64-hex Author ID
  (\texttt{.authorId}), which appears on the wire as
  \texttt{signerPubkey}. It has no location and no node ID ---
  authorship is not a place. Pass \texttt{persistAs} to keep a stable
  author across runs.
\end{itemize}

\textbf{The one rule:} every publish names its signer via
\texttt{\{\ signWith\ \}} --- an author identity, or the
\texttt{ANONYMOUS} sentinel for a deliberately unsigned post. There is
no default signer, and the node key never signs publishes.

\texttt{connect()} (below) mints both identities for you. To make the
author \emph{durable} across runs, pass
\texttt{author:\ \textquotesingle{}myapp:author\textquotesingle{}} --- a
\texttt{persistAs} key --- instead of the ephemeral default. The
underlying factories stay available for custom wiring.

\subsection{3. Topics are descriptors, not
strings}\label{topics-are-descriptors-not-strings}

A topic is a small descriptor
\texttt{\{\ region?,\ owner?,\ name,\ write?\ \}}; its 66-hex topic ID
is a hash of that descriptor. Because the ID is a pure function of the
fields, \textbf{anyone who knows the fields computes the same ID} --- no
registry, no coordination.

\begin{itemize}
\tightlist
\item
  \texttt{region} --- a real geographic cell. Accepts a name
  (\texttt{\textquotesingle{}useast\textquotesingle{}}), a hex string
  (\texttt{\textquotesingle{}0x89\textquotesingle{}}), or a number
  (\texttt{137}) --- all equivalent. Omit it and it defaults to the
  publisher's own node region; name it to share a topic across regions.
\item
  \texttt{owner} --- an Author ID, or absent.
\item
  \texttt{name} --- the human label
  (\texttt{\textquotesingle{}lobby\textquotesingle{}},
  \texttt{\textquotesingle{}profile\textquotesingle{}}).
\item
  \texttt{write} --- \texttt{\textquotesingle{}open\textquotesingle{}}
  or \texttt{\textquotesingle{}owner\textquotesingle{}}.
  \textbf{Defaults by owner:} no owner -\textgreater{}
  \texttt{\textquotesingle{}open\textquotesingle{}} (anyone publishes);
  an owner -\textgreater{}
  \texttt{\textquotesingle{}owner\textquotesingle{}} (only the owner key
  publishes, the safe default). With no owner, \texttt{write} is ignored
  --- you can't have an owner-only topic with no owner.
\end{itemize}

This Quick Start uses an \textbf{open} topic
\texttt{\{\ region,\ name\ \}}: anyone publishes, anyone subscribes.
(Note: the removed v2 model --- string topics, a \texttt{publisher}
argument, \texttt{sign:false},
\texttt{geoCellId}/\texttt{regionSynthPublisher} synthetic publishers
--- is gone. Topics are descriptors now.)

\texttt{deriveTopicId(descriptor)} returns the 66-hex ID --- the
shareable read handle. \texttt{sub} accepts \textbf{either} a descriptor
\textbf{or} that ID; \texttt{pub} always needs the descriptor (the
storing node recomputes the ID to enforce the write policy).

\subsection{4. Write the demo (one file)}\label{write-the-demo-one-file}

Save this as \texttt{index.js} --- it mirrors how
\texttt{apps/axona-minimal} wires a peer:

\begin{Shaded}
\begin{Highlighting}[]
\ImportTok{import}\NormalTok{ \{ connect}\OperatorTok{,}\NormalTok{ deriveTopicId}\OperatorTok{,}\NormalTok{ KERNEL\_VERSION \} }\ImportTok{from} \StringTok{\textquotesingle{}@axona/protocol\textquotesingle{}}\OperatorTok{;}

\KeywordTok{const}\NormalTok{ BRIDGE }\OperatorTok{=} \StringTok{\textquotesingle{}wss://testnet.axona.net\textquotesingle{}}\OperatorTok{;}           \CommentTok{// live testnet bridge (kernel 4.40.0)}
\KeywordTok{const}\NormalTok{ HERE   }\OperatorTok{=}\NormalTok{ \{ }\DataTypeTok{lat}\OperatorTok{:} \FloatTok{38.0}\OperatorTok{,} \DataTypeTok{lng}\OperatorTok{:} \OperatorTok{{-}}\FloatTok{77.0}\NormalTok{ \}}\OperatorTok{;}           \CommentTok{// your real location (us{-}east here)}
\KeywordTok{const}\NormalTok{ TOPIC  }\OperatorTok{=}\NormalTok{ \{ }\DataTypeTok{region}\OperatorTok{:} \StringTok{\textquotesingle{}useast\textquotesingle{}}\OperatorTok{,} \DataTypeTok{name}\OperatorTok{:} \StringTok{\textquotesingle{}quick{-}start{-}demo\textquotesingle{}}\NormalTok{ \}}\OperatorTok{;}   \CommentTok{// open topic}

\CommentTok{// 1–3. One call: mints both identities (connection + author), builds the}
\CommentTok{//      peer on the web transport, starts it, and waits for the mesh to warm.}
\BuiltInTok{console}\OperatorTok{.}\FunctionTok{log}\NormalTok{(}\StringTok{\textquotesingle{}kernel v\textquotesingle{}} \OperatorTok{+}\NormalTok{ KERNEL\_VERSION }\OperatorTok{+} \StringTok{\textquotesingle{} — connecting…\textquotesingle{}}\NormalTok{)}\OperatorTok{;}
\KeywordTok{const}\NormalTok{ \{ peer}\OperatorTok{,}\NormalTok{ author}\OperatorTok{,}\NormalTok{ status}\OperatorTok{,}\NormalTok{ disconnect \} }\OperatorTok{=} \ControlFlowTok{await} \FunctionTok{connect}\NormalTok{(\{}
  \DataTypeTok{bridge}\OperatorTok{:}\NormalTok{   BRIDGE}\OperatorTok{,}
  \DataTypeTok{location}\OperatorTok{:}\NormalTok{ HERE}\OperatorTok{,}
\NormalTok{\})}\OperatorTok{;}
\BuiltInTok{console}\OperatorTok{.}\FunctionTok{log}\NormalTok{(}\StringTok{\textquotesingle{}mesh \textquotesingle{}} \OperatorTok{+}\NormalTok{ (status}\OperatorTok{.}\AttributeTok{ready} \OperatorTok{?} \StringTok{\textquotesingle{}ready\textquotesingle{}} \OperatorTok{:} \StringTok{\textquotesingle{}not ready\textquotesingle{}}\NormalTok{) }\OperatorTok{+} \StringTok{\textquotesingle{} (\textquotesingle{}} \OperatorTok{+}\NormalTok{ status}\OperatorTok{.}\AttributeTok{peers} \OperatorTok{+} \StringTok{\textquotesingle{} peers, \textquotesingle{}} \OperatorTok{+}\NormalTok{ status}\OperatorTok{.}\AttributeTok{ms} \OperatorTok{+} \StringTok{\textquotesingle{}ms)\textquotesingle{}}\NormalTok{)}\OperatorTok{;}
\BuiltInTok{console}\OperatorTok{.}\FunctionTok{log}\NormalTok{(}\StringTok{\textquotesingle{}topic id:\textquotesingle{}}\OperatorTok{,} \ControlFlowTok{await} \FunctionTok{deriveTopicId}\NormalTok{(TOPIC))}\OperatorTok{;}

\CommentTok{// 4. Subscribe. The handler gets an envelope: \{ msgId, seq, ts, topic, message, signerPubkey? \}.}
\CommentTok{//    It fires for every delivery — including the ones we publish below — so you}
\CommentTok{//    watch the counter climb live in the console.}
\KeywordTok{const}\NormalTok{ sub }\OperatorTok{=} \ControlFlowTok{await}\NormalTok{ peer}\OperatorTok{.}\FunctionTok{sub}\NormalTok{(TOPIC}\OperatorTok{,}\NormalTok{ (env) }\KeywordTok{=\textgreater{}}\NormalTok{ \{}
  \BuiltInTok{console}\OperatorTok{.}\FunctionTok{log}\NormalTok{(}\StringTok{\textquotesingle{}[recv]\textquotesingle{}}\OperatorTok{,}\NormalTok{ env}\OperatorTok{.}\AttributeTok{message}\OperatorTok{,} \StringTok{\textquotesingle{}from\textquotesingle{}}\OperatorTok{,}\NormalTok{ (env}\OperatorTok{.}\AttributeTok{signerPubkey} \OperatorTok{||} \StringTok{\textquotesingle{}anon\textquotesingle{}}\NormalTok{)}\OperatorTok{.}\FunctionTok{slice}\NormalTok{(}\DecValTok{0}\OperatorTok{,} \DecValTok{12}\NormalTok{))}\OperatorTok{;}
\NormalTok{\}}\OperatorTok{,}\NormalTok{ \{ }\DataTypeTok{since}\OperatorTok{:} \StringTok{\textquotesingle{}all\textquotesingle{}}\NormalTok{ \})}\OperatorTok{;}

\CommentTok{// 5. Publish on a loop: one signed message per second, each carrying an}
\CommentTok{//    incrementing counter, so the system is visibly, continuously running.}
\KeywordTok{let}\NormalTok{ n }\OperatorTok{=} \DecValTok{0}\OperatorTok{;}
\KeywordTok{const}\NormalTok{ timer }\OperatorTok{=} \PreprocessorTok{setInterval}\NormalTok{(}\KeywordTok{async}\NormalTok{ () }\KeywordTok{=\textgreater{}}\NormalTok{ \{}
  \KeywordTok{const}\NormalTok{ msg }\OperatorTok{=} \VerbatimStringTok{\textasciigrave{}tick \#}\SpecialCharTok{$\{}\OperatorTok{++}\NormalTok{n}\SpecialCharTok{\}}\VerbatimStringTok{\textasciigrave{}}\OperatorTok{;}
  \KeywordTok{const}\NormalTok{ msgId }\OperatorTok{=} \ControlFlowTok{await}\NormalTok{ peer}\OperatorTok{.}\FunctionTok{pub}\NormalTok{(TOPIC}\OperatorTok{,}\NormalTok{ msg}\OperatorTok{,}\NormalTok{ \{ }\DataTypeTok{signWith}\OperatorTok{:}\NormalTok{ author \})}\OperatorTok{;}
  \BuiltInTok{console}\OperatorTok{.}\FunctionTok{log}\NormalTok{(}\StringTok{\textquotesingle{}[pub ]\textquotesingle{}}\OperatorTok{,}\NormalTok{ msg}\OperatorTok{,} \StringTok{\textquotesingle{}(msgId \textquotesingle{}} \OperatorTok{+}\NormalTok{ msgId}\OperatorTok{.}\FunctionTok{slice}\NormalTok{(}\DecValTok{0}\OperatorTok{,} \DecValTok{12}\NormalTok{) }\OperatorTok{+} \StringTok{\textquotesingle{}…)\textquotesingle{}}\NormalTok{)}\OperatorTok{;}
\NormalTok{\}}\OperatorTok{,} \DecValTok{1000}\NormalTok{)}\OperatorTok{;}

\CommentTok{// 6. Keep running until Ctrl+C, then leave cleanly.}
\BuiltInTok{process}\OperatorTok{.}\FunctionTok{on}\NormalTok{(}\StringTok{\textquotesingle{}SIGINT\textquotesingle{}}\OperatorTok{,} \KeywordTok{async}\NormalTok{ () }\KeywordTok{=\textgreater{}}\NormalTok{ \{}
  \PreprocessorTok{clearInterval}\NormalTok{(timer)}\OperatorTok{;}
  \ControlFlowTok{await}\NormalTok{ sub}\OperatorTok{.}\FunctionTok{stop}\NormalTok{()}\OperatorTok{;}
  \ControlFlowTok{await} \FunctionTok{disconnect}\NormalTok{()}\OperatorTok{;}          \CommentTok{// leave + stop, gracefully}
  \BuiltInTok{process}\OperatorTok{.}\FunctionTok{exit}\NormalTok{(}\DecValTok{0}\NormalTok{)}\OperatorTok{;}
\NormalTok{\})}\OperatorTok{;}
\end{Highlighting}
\end{Shaded}

\subsection{5. Run}\label{run}

\begin{Shaded}
\begin{Highlighting}[]
\ExtensionTok{node}\NormalTok{ index.js}
\end{Highlighting}
\end{Shaded}

You should see the counter climb, one tick per second, until you stop it
with \textbf{Ctrl+C}:

\begin{verbatim}
kernel v4.40.0 — connecting…
mesh ready (4 peers, 1200ms)
topic id: 89a1b2c3…
[pub ] tick #1 (msgId 8e9d4b1a30c2…)
[recv] tick #1 from 4f2a9b0c1d3e
[pub ] tick #2 (msgId 2c7f0a9b41d3…)
[recv] tick #2 from 4f2a9b0c1d3e
[pub ] tick #3 (msgId 91b3e5c8007a…)
[recv] tick #3 from 4f2a9b0c1d3e
…
\end{verbatim}

\textbf{That's it.} You just:

\begin{itemize}
\tightlist
\item
  minted an Ed25519 connection identity (66-hex node ID,
  region-prefixed) and a location-free author identity (64-hex Author ID
  = \texttt{signerPubkey}),
\item
  connected to the live bridge over the web transport and warmed up a
  routing mesh,
\item
  derived a deterministic topic ID from \texttt{\{\ region,\ name\ \}},
\item
  published a \textbf{signed} envelope to the topic's root axons and
  received it back through a
  \texttt{since:\ \textquotesingle{}all\textquotesingle{}} subscriber.
\end{itemize}

\subsection{What just happened}\label{what-just-happened}

\begin{verbatim}
   peer.pub({region,name}, msg, {signWith: author})        peer.sub({region,name}, handler, {since:'all'})
        |                                                       |
        v                                                       v
   deriveTopicId({region, name})  ----- same fields ----->  deriveTopicId({region, name})
        |   (identical 66-hex topic id on both sides)           |
        v                                                       v
   build signed envelope  -> route to topic's K-closest    subscribe-k to the same K-closest root axons
   root axons (publish-k)                                       |
        |                                                       v
        +--------------------- fan-out ----------------->  handler(env): { msgId, seq, ts, topic, message, signerPubkey }
\end{verbatim}

Both sides compute the identical topic ID because the ID is a pure hash
of the descriptor fields. That ID-matching is the rule you can't break
--- same \texttt{\{\ region,\ name\ \}}, same topic.

\subsection{Going further}\label{going-further}

{\def\LTcaptype{none} % do not increment counter
\begin{longtable}[]{@{}
  >{\raggedright\arraybackslash}p{(\linewidth - 2\tabcolsep) * \real{0.5000}}
  >{\raggedright\arraybackslash}p{(\linewidth - 2\tabcolsep) * \real{0.5000}}@{}}
\toprule\noalign{}
\begin{minipage}[b]{\linewidth}\raggedright
Want to\ldots{}
\end{minipage} & \begin{minipage}[b]{\linewidth}\raggedright
Do this
\end{minipage} \\
\midrule\noalign{}
\endhead
\bottomrule\noalign{}
\endlastfoot
Make authorship stable across runs &
\texttt{createAuthorIdentity(\{\ persistAs:\ \textquotesingle{}me\textquotesingle{}\ \})} \\
Publish anonymously &
\texttt{peer.pub(topic,\ msg,\ \{\ signWith:\ ANONYMOUS\ \})} (import
\texttt{ANONYMOUS}) \\
Own a feed only you can write &
\texttt{\{\ region,\ owner:\ me.authorId,\ name:\ \textquotesingle{}profile\textquotesingle{}\ \}}
(write defaults to
\texttt{\textquotesingle{}owner\textquotesingle{}}) \\
Share a read-only handle & \texttt{await\ deriveTopicId(descriptor)}
-\textgreater{} hand out the 66-hex ID;
\texttt{sub}/\texttt{pull}/\texttt{metrics} accept it \\
Run against production & set
\texttt{BRIDGE\ =\ \textquotesingle{}wss://bridge.axona.net\textquotesingle{}}
(same 4.x line; a separate network from testnet) \\
See the full mental model &
\href{Axona-Programmer-Guide-v4.38.0.md}{Programmer Guide} \\
Look up a specific symbol & \href{Axona-API-Reference-v4.40.0.md}{API
Reference} \\
\end{longtable}
}

\subsection{Troubleshooting}\label{troubleshooting}

\textbf{\texttt{Cannot\ find\ module\ \textquotesingle{}@axona/protocol\textquotesingle{}}}
--- make sure \texttt{package.json} has \texttt{"type":\ "module"} and
the install completed. Re-run \texttt{npm\ install}.

\textbf{\texttt{Cannot\ mix\ BigInt\ and\ other\ types}} --- only
applies to \emph{manual} peer assembly on kernels before 4.14;
\texttt{connect()} never hits it, and since 4.14 \texttt{NeuronNode}
accepts the hex id directly.

\textbf{\texttt{peer.pub:\ name\ a\ signer…}} --- \texttt{signWith} is
required. Pass an author identity, or
\texttt{\{\ signWith:\ ANONYMOUS\ \}} for an unsigned publish. There is
no default signer.

\textbf{Connects but nothing arrives} --- you almost certainly published
before the mesh warmed up. Subscribing/publishing on a cold synaptome
routes to too few root axons. \texttt{await\ peer.ready()} (step 3)
before calling \texttt{pub}/\texttt{sub} --- it resolves once the
routing mesh is warm. On the public bridge this takes a few seconds.

\textbf{Can't reach the bridge / connection hangs} --- confirm outbound
\texttt{wss://} (TLS WebSocket) to \texttt{testnet.axona.net} is allowed
by your network. There is no fallback to a local process; the demo needs
the bridge to find peers.

\textbf{\texttt{UPGRADE\_REQUIRED} close code (4426)} --- a wire/version
mismatch. The testnet bridge runs kernel 4.40.0 (wire 4.0); install
\texttt{github:axona-net/axona-protocol\#v4.40.0}. Production runs the
same wire-4 line (one release behind), so 4426 there means a genuinely
stale pin --- upgrade to the bridge's kernel version (shown at its
\texttt{/healthz}).

\end{document}
